\documentclass[11pt]{article}

\usepackage[T1]{fontenc}
\usepackage{amsmath,amssymb,amsthm,mathtools}
\usepackage[margin=0.82in]{geometry}
\usepackage{booktabs,longtable,array}

\newtheorem{lemma}{Lemma}
\newtheorem{proposition}[lemma]{Proposition}
\newtheorem{corollary}[lemma]{Corollary}

\newcommand{\W}{\mathcal W}
\newcommand{\Dlt}{\Delta}
\newcommand{\pimin}{\pi_-}
\newcommand{\pimax}{\pi_+}
\newcolumntype{L}[1]{>{\raggedright\arraybackslash}p{#1}}

\title{Connected High-Event Registry}
\author{}
\date{Analytic ledger replacement, Packet C}

\begin{document}
\maketitle

\begin{abstract}
This note replaces the exact-executable arithmetic used in the high
spherical-shell staircase by displayed rational witnesses.  It contains
the complete seventeen-face localization registry, the three algebraic
splits, the three dashed (impossible) cap cells, the five global checkpoint
payments, the seven conditional payments, and the corrected cap--74 path.
All comparisons are strict rational comparisons after the two elementary
bounds \(333/106<\pi<22/7\).  In particular, no decimal approximation,
interval package, or external ledger is a logical input to this packet.
\end{abstract}

\section{Definitions and strict proxy convention}

For \(0<\rho<1\), put
\begin{equation}
 z=z_\rho:=\frac{\pi}{1-\rho},\qquad
 x=x_\rho:=z_\rho^2,\qquad
 k_m(\rho):=\sqrt{x_\rho+m(m+1)}.
 \label{eq:def-z-k}
\end{equation}
The three-dimensional Weyl term is
\begin{equation}
 \W(\rho,K):=\frac{2}{9\pi}(1-\rho^3)K^3.
 \label{eq:def-W}
\end{equation}
At \(K=k_m(\rho)\), the product part of the strict channel proxy for the
mode \((\ell,n)\), \(n\geq2\), is
\[
 n^2x+\ell(\ell+1)<x+m(m+1).
\]
It is therefore governed by
\begin{equation}
 \Dlt_{m;n,\ell}(\rho)
 :=m(m+1)-\ell(\ell+1)-(n^2-1)x_\rho.
 \label{eq:delta}
\end{equation}
The channel is product-possible only when \(\Dlt_{m;n,\ell}>0\).
Equality is excluded because the spectral proxy is strict.  Since
\begin{equation}
 x_\rho'=\frac{2\pi^2}{(1-\rho)^3}>0,
 \qquad
 \Dlt_{m;n,\ell}'=-(n^2-1)x_\rho'<0,
 \label{eq:delta-monotone}
\end{equation}
a positive witness at \(r\) controls the side \(\rho\leq r\), while a
negative witness controls the side \(\rho\geq r\).

We write \(H\) for the largest first-mode angular index retained in a cap
cell and \(h\) for the largest second-mode angular index.  A complete block
through \(H\) has multiplicity \((H+1)^2\), because
\(\sum_{\ell=0}^{H}(2\ell+1)=(H+1)^2\).

\section{The only transcendental enclosure}

\begin{lemma}[Rational enclosure of \(\pi\)]
\label{lem:pi}
One has
\begin{equation}
 \pimin:=\frac{333}{106}<\pi<\frac{22}{7}=: \pimax.
 \label{eq:pi-bounds}
\end{equation}
\end{lemma}

\begin{proof}
For \(0<u<1\), the alternating expansion of \(\arctan u\) lies between
each pair of consecutive partial sums.  The tangent addition formula gives
\[
 \tan\!\left(4\arctan\frac15-\arctan\frac1{239}\right)=1;
\]
indeed, \(\tan(4\arctan(1/5))=120/119\), so both the numerator and
denominator in the last tangent quotient are \(28561/(119\cdot239)\).
The angle lies in \((0,\pi/2)\), and hence Machin's identity follows:
\[
 \pi=16\arctan\frac15-4\arctan\frac1{239}.
\]
Using an even partial sum for the positive term and an upper partial sum
for the subtracted term gives
\begin{align*}
 \pi
 &>16\left(\frac15-\frac1{3\cdot5^3}
       +\frac1{5\cdot5^5}-\frac1{7\cdot5^7}\right)-\frac4{239},\\
 \left[\text{right side}\right]-\frac{333}{106}
 &=\frac{3418213}{41563593750}>0.
\end{align*}
Likewise,
\begin{align*}
 \pi
 &<16\left(\frac15-\frac1{3\cdot5^3}
       +\frac1{5\cdot5^5}\right)
       -4\left(\frac1{239}-\frac1{3\cdot239^3}\right),\\
 \frac{22}{7}-\left[\text{right side}\right]
 &=\frac{1845738322}{1493178640625}>0.
\end{align*}
This proves \eqref{eq:pi-bounds}.
\end{proof}

\section{The seventeen rational localization faces}

Let
\[
 \rho_c:=\frac1{1+2\pi}.
\]
Because \(\pi>3\), \(\rho_c<1/7\).  The seventeen registry ratios are
strictly ordered as follows:
\begin{equation}
\begin{split}
 \frac4{25}<\frac15<\frac14<\frac3{10}<\frac13<\frac7{20}
 <\frac38<\frac25<\frac5{12}<\frac37<\frac12\\
 <\frac8{15}<\frac{107}{200}<\frac35<\frac{16}{25}
 <\frac{13}{20}<\frac23.
\end{split}
\label{eq:ratio-order}
\end{equation}
For a completely explicit distinctness check, the sixteen consecutive
gaps in \eqref{eq:ratio-order} are
\begin{equation}
 \frac1{25},\frac1{20},\frac1{20},\frac1{30},\frac1{60},
 \frac1{40},\frac1{40},\frac1{60},\frac1{84},\frac1{14},
 \frac1{30},\frac1{600},\frac{13}{200},\frac1{25},
 \frac1{100},\frac1{60}.
 \label{eq:ratio-gaps}
\end{equation}
Moreover
\[
 \frac4{25}-\frac17=\frac3{175}>0,
 \qquad
 \frac78-\frac23=\frac5{24}>0.
\]
Thus all seventeen ratios lie in the high strip
\(\rho_c<\rho<7/8\).

For a positive row below, we use
\[
 \Dlt_{m;n,\ell}(r)>
 m(m+1)-\ell(\ell+1)
 -(n^2-1)\left(\frac{\pimax}{1-r}\right)^2.
\]
For a negative row, we use
\[
 -\Dlt_{m;n,\ell}(r)>
 (n^2-1)\left(\frac{\pimin}{1-r}\right)^2
 -m(m+1)+\ell(\ell+1).
\]
The resulting exact witnesses are the following.  ``Retain'' means only
that the product wall does not permit deletion of the indicated block;
the independent Bessel wall may still remove it.

\begingroup
\small
\renewcommand{\arraystretch}{1.16}
\begin{longtable}{@{}c c c L{3.15cm} L{5.35cm}@{}}
\toprule
no.&\(r\)&\((m;n,\ell)\)&exact sign witness&one-sided cap/event consequence\\
\midrule
\endfirsthead
\toprule
no.&\(r\)&\((m;n,\ell)\)&exact sign witness&one-sided cap/event consequence\\
\midrule
\endhead
1&\(1/5\)&\((11;3,1)\)&
 \(\Dlt>320/49\)&retain the third block \(0{:}1\) for \(\rho\leq1/5\)\\
2&\(3/10\)&\((8;2,3)\)&
 \(-\Dlt>58215/137641\)&at \(k_8\), \(h\leq2\) for \(\rho\geq3/10\)\\
3&\(1/3\)&\((8;2,2)\)&
 \(-\Dlt>27699/44944\)&at \(k_8\), \(h\leq1\) for \(\rho\geq1/3\)\\
4&\(3/8\)&\((9;2,3)\)&
 \(\Dlt>2622/1225\)&retain the \(h=3\) column below this face\\
5&\(2/5\)&\((9;2,2)\)&
 \(\Dlt>248/147\)&retain the \(h=2\) column below this face\\
6&\(5/12\)&\((9;2,1)\)&
 \(\Dlt>2200/2401\)&retain the \(h=1\) column below this face\\
7&\(3/7\)&\((9;2,0)\)&
 \(-\Dlt>120843/179776\)&no second mode at \(k_9\) on the upper side\\
8&\(1/2\)&\((10;2,0)\)&
 \(-\Dlt>23677/2809\)&no second mode at \(k_{10}\) on the upper side\\
9&\(107/200\)&\((11;2,0)\)&
 \(-\Dlt>13302732/2699449\)&no second mode at \(k_{11}\) on the upper side\\
10&\(8/15\)&\((11;2,0)\)&
 \(-\Dlt>2175627/550564\)&the clean stated cutoff: a second mode forces \(\rho<8/15\)\\
11&\(3/5\)&\((8;2,0)\)&
 \(-\Dlt>5080707/44944\)&select the no-second \(k_8\) cap/payment side\\
12&\(16/25\)&\((9;2,0)\)&
 \(-\Dlt>1555635/11236\)&select the no-second \(k_9\) cap/payment side\\
13&\(13/20\)&\((10;2,0)\)&
 \(-\Dlt>18126190/137641\)&select the no-second \(k_{10}\) cap/payment side\\
14&\(2/3\)&\((11;2,0)\)&
 \(-\Dlt>1510851/11236\)&select the no-second \(k_{11}\) cap/payment side\\
15&\(4/25\)&\((10;3,0)\)&
 \(-\Dlt>260740/137641\)&no third mode at \(k_{10}\) on the upper side\\
16&\(1/4\)&\((11;3,0)\)&
 \(-\Dlt>23484/2809\)&a third mode at \(k_{11}\) forces \(\rho<1/4\)\\
17&\(7/20\)&\((9;2,4)\)&
 \(-\Dlt>36230/474721\)&the \(H=6,h=4\) cap--74 cell can occur only below this face\\
\bottomrule
\end{longtable}
\endgroup

Every entry is a displayed positive rational number.  Together with the
monotonicity \eqref{eq:delta-monotone}, the table proves the asserted side
on an interval, not merely at a sampled point.

\section{Three algebraic splits and the five \(k_{11}\) implications}

At \(\rho=\rho_c\), \(x=(\pi+1/2)^2\), while at \(\rho=7/8\),
\(x=64\pi^2\).  Lemma~\ref{lem:pi} gives the exact chain
\begin{equation}
 (\pi+\tfrac12)^2
 <\left(\frac{51}{14}\right)^2
 =\frac{2601}{196}<16<\frac{68}{3}<34<64\cdot3^2<64\pi^2.
 \label{eq:split-strip}
\end{equation}
Consequently each of \(x=16,68/3,34\) is attained at a unique point in
the high strip.  The relevant channel arithmetic is
\begin{equation}
\begin{array}{c|c|c}
x&\text{channel}&\text{exact value of }\Dlt\\ \hline
16&(m;n,\ell)=(11;3,1)&132-2-8\cdot16=2>0,\\
68/3&(10;2,6)&110-42-3(68/3)=0,\\
34&(11;2,5)&132-30-3\cdot34=0.
\end{array}
\label{eq:algebraic-splits}
\end{equation}
Thus the first split lies on the possible side.  The latter two are exact
channel walls, and their new modes are excluded by the strict convention.

The two localization consequences at \(k_{11}\) may also be read without
the table:
\begin{align}
 x_{8/15}-44
 &>\left(\frac{\pimin}{1-8/15}\right)^2-44
 =\frac{725209}{550564}>0,\label{eq:k11-second}\\
 x_{1/4}-\frac{33}{2}
 &>\left(\frac{\pimin}{1-1/4}\right)^2-\frac{33}{2}
 =\frac{5871}{5618}>0.\label{eq:k11-third}
\end{align}
Hence a second mode at \(k_{11}\) forces \(\rho<8/15\), and a third mode
forces \(\rho<1/4\).  Finally, the first-mode ball walls give
\begin{align}
 \left(\frac{69}{5}\right)^2-132
 &=\frac{1461}{25}
 =44+\frac{361}{25}
 =\frac{33}{2}+\frac{2097}{50},\label{eq:H10}\\
 \left(\frac{64}{5}\right)^2-132
 &=\frac{796}{25}
 =\frac{33}{2}+\frac{767}{50}.\label{eq:H9}
\end{align}
If \(H=10\) is present at \(k_{11}\), then \(x\) is larger than both the
second threshold \(44\) and the third threshold \(33/2\); hence every
second and third mode is excluded.  If \(H=9\) is present, every third
mode is excluded.  Equations \eqref{eq:k11-second}--\eqref{eq:H9} are the
five formerly separate \(k_{11}\) localization implications.

\section{The three impossible cap cells}

If the first mode \(H\) is counted at \(k_m\), its strict ball wall
\(\beta_{H,1}<k_m\) forces
\(x>\beta_{H,1}^2-m(m+1)\).  If the second mode \(h\) is counted, its
product wall forces
\(x<(m(m+1)-h(h+1))/3\).  The three dashed cells are therefore empty by
the following exact, incompatible bounds:
\begin{equation}
\begin{array}{c|c|c|c}
\text{cell}&\text{required lower bound}&\text{required upper bound}
&\text{gap}\\ \hline
k_8: H=6,\ h\ge2&100-72=28&(72-6)/3=22&6\\
k_9: H=8,\ h\ge0&(71/6)^2-90=1801/36&90/3=30&721/36\\
k_9: H=7,\ h\ge3&(23/2)^2-90=169/4&(90-12)/3=26&65/4.
\end{array}
\label{eq:impossible-cells}
\end{equation}
No limiting equality can rescue a cell: both necessary modal inequalities
are strict.

\section{Checkpoint monotonicity and five global payments}

Put \(M=m(m+1)\) and
\(\W_m(\rho):=\W(\rho,k_m(\rho))\).  Logarithmic differentiation shows
that \(\W_m'\) has the sign of
\begin{equation}
 \pi^2(1+\rho)-M\rho^2(1-\rho)^2.
 \label{eq:Wm-derivative}
\end{equation}
For \(m\leq11\),
\[
 \pi^2(1+\rho)-M\rho^2(1-\rho)^2
 >9-\frac{132}{16}=\frac34>0.
\]
Thus every \(\W_m\), \(m\leq11\), is strictly increasing.  This single
argument replaces both executable derivative-bound predicates.

At \(\rho_c\), one has \(z=\pi+1/2>193/53\),
\(\rho_c<1/7\), and \(1/\pi>7/22\).  Consequently
\begin{equation}
 \W(\rho_c,K)>
 \frac{2}{9}\frac7{22}\left(1-\frac1{7^3}\right)K^3
 =\frac{38}{539}K^3.
 \label{eq:global-coarse}
\end{equation}
For each row in the next table, the middle column proves
\(k_m(\rho_c)>t_m\), and the last column proves
\(\W_m(\rho_c)>C_m\).

\begin{center}
\small
\renewcommand{\arraystretch}{1.18}
\begin{tabular}{@{}c c c c c@{}}
\toprule
\(m\)&cap \(C_m\)&\(t_m\)&
\((193/53)^2+M-t_m^2\)&\((38/539)t_m^3-C_m\)\\
\midrule
7&40&\(208/25\)&\(67049/1755625\)&\(5083656/8421875\)\\
8&53&\(923/100\)&\(1901439/28090000\)&\(656778873/269500000\)\\
9&73&\(254/25\)&\(61431/1755625\)&\(7911557/8421875\)\\
10&89&\(111/10\)&\(14211/280900\)&\(1999489/269500\)\\
11&121&\(241/20\)&\(65271/1123600\)&\(5076899/2156000\)\\
\bottomrule
\end{tabular}
\end{center}
Every last-column entry is positive.  Monotonicity therefore extends each
global payment from \(\rho_c\) to the complete high strip.

\section{Seven conditional first-block payments}

We now give a general reduction which avoids any numerical enclosure of a
radical.  Suppose the presence of first mode \(H\) below the checkpoint
\(k_m\) forces \(K>t\).  Set
\begin{equation}
 a^2:=t^2-m(m+1),\qquad
 \rho_0:=1-\frac{\pi}{a}.
 \label{eq:conditional-rho0}
\end{equation}
All seven rows below have \(a>\pi\): the smallest listed radicand is
\[
 \frac{89}{4},\qquad
 \frac{89}{4}-\left(\frac{22}{7}\right)^2
 =\frac{2425}{196}>0.
\]
If the mode is present for some
\(K\leq k_m(\rho)\), then \(k_m(\rho)>t\), hence
\(x_\rho>a^2\) and \(\rho>\rho_0\).  By the strict increase of
\(\W_m\),
\[
 \W_m(\rho)>\W_m(\rho_0)=\W(\rho_0,t).
\]
Choose a rational \(r>\rho_0\).  Since \(\W(\rho,t)\) decreases with
\(\rho\) at fixed \(t\), Lemma~\ref{lem:pi} yields
\begin{equation}
 \W(\rho_0,t)>\W(r,t)>
 \frac7{99}(1-r^3)t^3.
 \label{eq:conditional-coarse}
\end{equation}
The condition \(r>\rho_0\) follows, without taking a square root, from
\begin{equation}
 L:=\pimin^2-(1-r)^2a^2>0.
 \label{eq:location-witness}
\end{equation}
Here is the complete calculation.  The column \(P\) is the reserve in
\eqref{eq:conditional-coarse} above the stated cap.

\begingroup
\small
\renewcommand{\arraystretch}{1.18}
\begin{longtable}{@{}c c c c c c c c@{}}
\toprule
\(m\)&\(H\)&\(t\)&cap&\(a^2\)&\(r\)&\(L\)&
\(P=\frac7{99}(1-r^3)t^3-\mathrm{cap}\)\\
\midrule
\endfirsthead
\toprule
\(m\)&\(H\)&\(t\)&cap&\(a^2\)&\(r\)&\(L\)&\(P\)\\
\midrule
\endhead
7&6&\(10\)&53&44&\(8/15\)&\(725209/2528100\)&\(18659/2673\)\\
8&7&\(23/2\)&80&\(241/4\)&\(3/5\)&\(64349/280900\)&\(213281/49500\)\\
9&8&\(71/6\)&81&\(1801/36\)&\(3/5\)&\(4713989/2528100\)&\(14506973/1336500\)\\
10&7&\(23/2\)&100&\(89/4\)&\(7/20\)&\(2105431/4494400\)&\(18539033/6336000\)\\
10&8&\(71/6\)&97&\(1081/36\)&\(3/7\)&\(317585/4955076\)&\(2865431/261954\)\\
10&9&\(64/5\)&100&\(1346/25\)&\(3/5\)&\(8811001/7022500\)&\(25143284/1546875\)\\
11&9&\(64/5\)&125&\(796/25\)&\(1/2\)&\(536261/280900\)&\(58757/12375\)\\
\bottomrule
\end{longtable}
\endgroup
All \(L\)'s and \(P\)'s are positive.  Each row therefore proves both
the localization of its conditional branch and its Weyl payment.

\section{The connected cap--74 path}

The seventeenth localization row is tied to the exceptional
\(k_9,H=6,h=4\) cap as follows.  Angular propagation from the strict
\((3,2)\) Bessel wall gives
\begin{equation}
 j_{9/2,2}^2>\left(\frac{103}{10}\right)^2+20-12
 =\frac{11409}{100}=108+\frac{609}{100}.
 \label{eq:cap74-ball}
\end{equation}
Accordingly the strict proxy uses
\(\beta_{4,2}^2=11409/100\); the strict inequality belongs to the actual
Bessel zero \(j_{9/2,2}\), not to the proxy constant.
At \(k_9\), the product wall for \((4,2)\) is
\[
 3x<90-20=70,
 \qquad\text{i.e.}\qquad x<\frac{70}{3}.
\]
The ball wall \eqref{eq:cap74-ball}, on the other hand, forces
\[
 x>\frac{11409}{100}-90=\frac{2409}{100},
 \qquad
 \frac{2409}{100}-\frac{70}{3}=\frac{227}{300}>0.
\]
Thus the fully sharpened cap--74 cell is empty.

For completeness, even the deliberately conservative cell is paid.  At
\(\rho=7/20\),
\begin{equation}
 x_{7/20}-\frac{70}{3}
 >\frac{400}{169}\left(\frac{333}{106}\right)^2-\frac{70}{3}
 =\frac{36230}{1424163}>0,
 \label{eq:cap74-rho}
\end{equation}
so a counted \((4,2)\) mode forces \(\rho<7/20\).  Also
\begin{equation}
 108-\left(\frac{207}{20}\right)^2=\frac{351}{400}>0,
 \label{eq:cap74-K}
\end{equation}
so its ball wall forces \(K>207/20\).  Therefore
\begin{align}
 \W(\rho,K)
 &>\frac7{99}\left(1-\left(\frac7{20}\right)^3\right)
       \left(\frac{207}{20}\right)^3\notag\\
 &=\frac{52823261673}{704000000}
 =74+\frac{727261673}{704000000}>74.
 \label{eq:cap74-payment}
\end{align}
The cap itself is exactly
\((6+1)^2+(4+1)^2=49+25=74\).  Equations
\eqref{eq:cap74-ball}--\eqref{eq:cap74-payment} cover both the actual
incompatibility and every row of the conservative cap--74 fallback.

\section{Exact row map and scope audit}

The following map uses the ordering of the two functions in the original
exact registry.  To keep the table readable, \(L\) denotes
\texttt{verify\_localizations}, and \(C\) denotes
\texttt{verify\_cap\_tables\_and\_payments}.  The map lets a referee match
every replaced predicate to a displayed hand calculation.

\begin{center}
\small
\renewcommand{\arraystretch}{1.16}
\begin{tabular}{@{}c c L{12.1cm}@{}}
\toprule
registry family&row numbers&analytic replacement here\\
\midrule
\(L\)&1--2&distinctness and connectedness:
\eqref{eq:ratio-order}--\eqref{eq:ratio-gaps} and the seventeen-row table\\
\(L\)&3--19&all ratios in the high strip, by
\(\rho_c<1/7<4/25\) and \(2/3<7/8\)\\
\(L\)&20--36&the seventeen exact sign witnesses in
the long table\\
\(L\)&37--42&split location
\eqref{eq:split-strip} and side/equality arithmetic
\eqref{eq:algebraic-splits}\\
\(L\)&43--47&the five \(k_{11}\) implications
\eqref{eq:k11-second}--\eqref{eq:H9}\\
\(C\)&44--55&the twelve non-bookkeeping cap--74
predicates, \eqref{eq:cap74-ball}--\eqref{eq:cap74-payment};
row \(C56\), the block identity \(49+25=74\), is allocated to Packet~A\\
\(C\)&57--59&the three impossible
cells, \eqref{eq:impossible-cells}\\
\(C\)&78--79&the general checkpoint
monotonicity proof \eqref{eq:Wm-derivative}\\
\(C\)&80--84&the five global-payment
rows following \eqref{eq:global-coarse}\\
\(C\)&85--98&the seven localization
and seven conditional-payment rows in the final long table\\
\bottomrule
\end{tabular}
\end{center}

The count reconciliation is
\[
\begin{array}{c|c|c}
\text{subfamily}&\text{connected mathematical objects}
&\text{registry predicates}\\ \hline
\text{registry metadata}&1&2\\
\text{high-strip membership}&17&17\\
\text{localization signs}&17&17\\
\text{algebraic splits}&3&6\\
\text{special }k_{11}\text{ implications}&5&5\\
\text{cap--74 path}&1&12\\
\text{impossible cells}&3&3\\
\text{checkpoint monotonicity}&1&2\\
\text{global payments}&5&5\\
\text{conditional payments}&7&14\\ \hline
\text{total}&60&83.
\end{array}
\]
In particular, the five headline Packet--C subfamilies contain exactly
\[
 17+3+3+5+7=35
\]
connected objects (localization signs, algebraic splits, impossible cells,
global payments, and conditional payments).  Because each algebraic split
has a location and a side predicate, and each conditional payment has a
localization and a reserve predicate, these headline objects replace
\(17+6+3+5+14=45\) registry predicates.  The remaining 38 predicates in
the total are the explicitly listed auxiliary closure rows.

Thus this packet replaces all 47 predicates of \(L\) and 36 predicates of
\(C\).  The elementary cap addition rows \(1\)--\(43\), row \(56\), and
rows \(60\)--\(77\) of the latter family are
outside the declared Packet~C scope; they are instances of the general
odd-sum/block-cap identity and are not claimed here.  No row in the
declared scope is unreconstructed.

\end{document}
